% "World of Adventure" campaign guide - Crafting
% Written by Christopher Thomas.

\chapter{Crafting}
\label{sect-crafting}

%
\section{Making Things}

To craft an item, you make weekly Craft checks against the item's crafting
DC until it's crafted:

\begin{itemize}
%
\item Successful checks add
\mbox{\bfseries 3 $\times$ check result $\times$ income multiplier sp}
towards the value of the object being crafted.
%
\item ``Income multiplier'' is usually x1 for crafting purposes. At the DM's
discretion, it might be x2 for jewellery and clockwork, x$\frac{1}{2}$ for
things like making bricks, and so forth.
%
\item Successful checks also consume
\mbox{\bfseries 1.5 $\times$ check result $\times$ income multiplier sp}
in materials and other supplies, but it's easier just to say ``producing an
item costs half its price in materials'' and deduct that as a lump sum.
%
\item Checks that fail by 5 or less make no progress but consume no
additional cost; you wasted time but avoided wasting materials or money.
%
\item Checks that fail by 6-10 lose
\mbox{\bfseries 1 $\times$ craft DC $\times$ income multiplier sp}
worth of materials or supplies. Checks that fail by more than 10 lose twice
this amount.
%
\end{itemize}

The net result of this is that if you're cranking out items as fast as you
can make them, always pass your skill checks, and always have a buyer,
you'll make about 1.5x as much money as you would if making income rolls.
If you meet the target DC with a die roll of 9 or more, you'll succeed
about two thirds of the time, very rarely waste materials, and average out
to about the amount predicted by income rolls. The income rolls abstract
the ``occasionally flub rolls'' and ``have to find buyers'' parts of the
crafting process.

%
\subsection{Making Things for Yourself}

Making items for yourself might save time and money, or might cost time
and money. You don't have to wait for a specialist to become available, but
how well it works depends on the crafting DC:

\begin{itemize}
\item If a die roll of 5 meets the target DC, you'll proceed at
\textbf{90\% speed} on average.
\item If a die roll of 7 meets the target DC, you'll proceed at
\textbf{80\% speed} on average.
\item If a die roll of 9 meets the target DC, you'll proceed at
\textbf{two thirds speed} on average and usually won't waste material.
This is the break-even point for making vs. buying.
\item If you need a die roll higher than 9 to meet the target DC, you'll
be spending more time and money than if you'd just commissioned a vendor to
make the item.
\end{itemize}

It can absolutely still be worth making an item for personal reasons even
if it takes more time and money than buying one.

%
\subsection{Making Magic Items}

A permanent magic item's physical vessel is a masterwork object. This is
made using an appropriate Craft skill (such as jewellery smithing or
armour smithing), and making the object progresses as described above.

To create an item suitable for enchantment, the person crafting the object
must have a number of ranks in Spellcraft equal to the number of
``enchantability'' masterworking improvements.

For any type of magic item (expendable or permanent), \textit{after} any
needed physical vessel is crafted, the enchantment is applied by making
crafting checks using the Spellcraft skill:

\begin{itemize}
\item Successful checks add \mbox{\bfseries 20 $\times$ check result sp}
per week. There is no ``income multiplier'' (it's already included).
\item DC is typically \mbox{\bfseries 15 + spell level} for scrolls,
potions, and similar one-shot items, and \mbox{\bfseries 20 + spell level}
or \mbox{\bfseries 20 + bonus} for other items.
\end{itemize}

This progresses much more quickly than most other forms of crafting. This
is because most of the time spent by a professional wizard isn't spent
crafting items - it's spent finding customers willing to commission items.

\fixme{This is no longer true for the ``magic is even cheaper'' scenario.
Consider a sliding scale of 6x for low cost to 20x for high cost.}

Magic item costs are described in detail in Section \ref{sect-magic}.
%

%
\section{Crafting DCs}

For most items, crafting DC goes up as the logarithm of the item's base
price (before masterworking):

\begin{tabular}{lc}\hline
\textbf{Base Price} & \textbf{Craft DC} \\ \hline
up to 1 sp & 11 \\
up to 3 sp & 13 \\
up to 10 sp & 15 \\
up to 30 sp & 17 \\
up to 100 sp & 19 \\
up to 300 sp & 21 \\
up to 1000 sp & 23 \\
\multicolumn{2}{c}{\ldots} \\ \hline
\end{tabular}

Exact crafting DCs are subject to DM discretion.

The difficulty classes that various types of people and organizations can
hit are roughly as follows:

\begin{itemize}
\item A decent L1 novice can hit DC~15 reliably, DC~20 sometimes, and DC~25
by fluke.
\item A decent late-career professional can hit DC~25 reliably, DC~30
sometimes, and DC~35 by fluke.
\item The very best solo practitioners (maybe with one helper) can hit DC~35
reliably, DC~40 sometimes, and DC~45 by fluke.
\item The very best organizations can hit DC~45 reliably, DC~50 sometimes,
and DC~55 by fluke.
\end{itemize}

Having people help you is described in Section \ref{sect-money}. A group
of helpers, in addition to contributing to crafting \textit{progress}, also
gives a bonus to the leader's crafting roll. The end result is that an
organization run by decent late-career professionals can match the DCs of
the best solo practitioners, while making crafting progress much more
rapidly (at the cost of organizational overhead).

%
% This is the end of the file.
